\section*{Hoofdstuk \ref{ch:g4:what-a-fire-needs} --- Branden: wat een vuur nodig heeft}

\begin{solution}{exo:g4:what-a-fire-needs:1}
\omterm{def:g4:what-a-fire-needs:fuel}{Brandstof}, zuurstof (uit de \omterm{def:g4:air-a-mixture-of-gases:air}{lucht}) en warmte (om het vuur aan te
steken).
\end{solution}

\begin{solution}{exo:g4:what-a-fire-needs:2}
Hout, kaarsvet, benzine en papier zijn \omterm{def:g4:what-a-fire-needs:fuel}{brandstoffen}. Zand, glas en water
\omterm{def:g4:what-a-fire-needs:burning}{branden} niet.
\end{solution}

\begin{solution}{exo:g4:what-a-fire-needs:3}
De zuurstof: de deken zorgt dat de \omterm{def:g4:air-a-mixture-of-gases:air}{lucht} het vuur niet bereikt.
\end{solution}

\begin{solution}{exo:g4:what-a-fire-needs:4}
De warmte: het water koelt het brandende hout zo sterk af dat het niet
meer kan \omterm{def:g4:what-a-fire-needs:burning}{branden}.
\end{solution}

\begin{solution}{exo:g4:what-a-fire-needs:5}
Drie van deze: as, rook, koolstofdioxide, waterdamp.
\end{solution}

\begin{solution}{exo:g4:what-a-fire-needs:6}
Blazen brengt meer \omterm{def:g4:air-a-mixture-of-gases:air}{lucht}, en dus meer zuurstof, bij de hete \omterm{def:g4:what-a-fire-needs:fuel}{brandstof}, en
die gaat weer \omterm{def:g4:what-a-fire-needs:burning}{branden}.
\end{solution}

\begin{solution}{exo:g4:what-a-fire-needs:7}
Water: het brandende kaarsvet maakt waterdamp, die op het koude glas
druppeltjes vormt.
\end{solution}

\begin{solution}{exo:g4:what-a-fire-needs:8}
Naar buiten gaan, laag onder de rook blijven en de deuren achter je
dichtdoen, en dan van buitenaf de brandweer bellen.
\end{solution}

\begin{solution}{exo:g4:what-a-fire-needs:9}
De brandstofzijde: met het gas dicht heeft het vuur niets meer om te
\omterm{def:g4:what-a-fire-needs:burning}{verbranden}.
\end{solution}

\begin{solution}{exo:g4:what-a-fire-needs:10}
De \omterm{def:g4:what-a-fire-needs:fuel}{brandstof}: in de strook staan geen bomen meer, dus als het vuur daar
aankomt, heeft het geen \omterm{def:g4:what-a-fire-needs:fuel}{brandstof} meer.
\end{solution}

\begin{solution}{exo:g4:what-a-fire-needs:11}
Nee. Het meeste hout is samen met zuurstof uit de \omterm{def:g4:air-a-mixture-of-gases:air}{lucht} nieuwe stoffen
geworden die gassen zijn --- koolstofdioxide, waterdamp en rook --- en
die zijn de \omterm{def:g4:air-a-mixture-of-gases:air}{lucht} in gegaan. Alleen de as blijft achter.
\end{solution}
